% ---------------------------------------------------------------------
%  5. Part Two: the science (handoff negative result is the headline)
% ---------------------------------------------------------------------
\section{Part Two: which cooperation actually pays}
\label{sec:parttwo}

Varying only the roles gives the per-problem picture in Table~\ref{tab:perproblem}.
I read it as exploratory evidence ($n=9$, single runs). A partial three-seed
replication of the integrated agent reproduced the reliable core and confirmed
that the fifth win is stochastic (\S\ref{sec:results}); the same-budget router
baseline and the full seeded matrix were not completed, so these cells stand as
single-run patterns rather than pass-rates.

\begin{table}[t]
\centering
\caption{Per-problem substantive outcomes on \Sdev{} (a dot is a substantive
pass). The \texttt{putnam\_2018\_a1} row is shown separately: its GPT-side
comparator ``passes'' are the answer-set tautology of \S\ref{sec:mechanism} and are
not substantive, so they do not enter the totals. All nine \Sdev{} problems are
listed; an empty row is an all-condition failure.}
\label{tab:perproblem}
\begin{tabular}{@{}lcccccc@{}}
\toprule
Problem & \SQ & \SG & \RQ & \RG & \HQG & \HGQ \\
\midrule
p01\_linear         & $\bullet$ & $\bullet$ & $\bullet$ & $\bullet$ & $\bullet$ & $\bullet$ \\
p03\_sq\_ge\_two\_ab & $\bullet$ & $\bullet$ & $\bullet$ & $\bullet$ & $\bullet$ & $\bullet$ \\
p05\_gcd\_mersenne   &           & $\bullet$ &           & $\bullet$ &           &           \\
p06\_pow\_mod        & $\bullet$ &           & $\bullet$ &           & $\bullet$ &           \\
p10\_factorial\_pow  &           &           & $\bullet$ & $\bullet$ &           &           \\
p09\_imo1964         &           &           &           &           &           &           \\
rmo\_2000\_2         &           &           &           &           &           &           \\
rmo\_2000\_3         &           &           &           &           &           &           \\
\midrule
\textbf{Substantive /9} & \textbf{3} & \textbf{3} & \textbf{4} & \textbf{4} & \textbf{3} & \textbf{2} \\
\midrule
putnam\_2018\_a1 (comparator only) & & $\circ$ & & $\circ$ & $\circ$ & $\circ$ \\
\bottomrule
\end{tabular}
\end{table}

\subsection{Finding 1 --- the models are complementary}
Their easy wins overlap; their harder wins do not. Among the model-driven arms,
\texttt{p06} is solved only on the Qwen side and \texttt{p05} only on the GPT side.
(In the integrated agent \texttt{p05} is closed even more cheaply, by the tactic
sweep, before any model runs; the complementarity here is a statement about the
two \emph{models}, read from the model-only arms.) Because the failures are
complementary, the derived Solo Union reaches 4/9 with no cross-model protocol,
and the cheapest correct move is to \emph{route} rather than merge.

\subsection{Finding 2 --- same-model repair supplies the margin}
Feeding the failed proof and the exact Lean messages back to the \emph{same} model
is what lifts the score above the union: both repair arms unlock
\texttt{p10\_factorial\_pow}, which no solo or handoff arm solves. That one problem
is the difference between the 4/9 union and the 5/9 portfolio. Compiler-in-the-loop
repair is the real lever --- and note it requires only \emph{one} model.

\subsection{Finding 3 (the headline, and it is negative) --- cross-model handoff
does not pay}
The intuitive hypothesis of the whole exercise is that two models cooperating ---
one repairing the other --- beat one model alone. On these two models, \emph{it
does not.} \HQG{} (Qwen proposes, GPT repairs) scores 3/9 and adds no problem that
same-model repair did not already solve. \HGQ{} (GPT proposes, Qwen repairs) is the
\emph{worst} arm at 2/9, and it \emph{loses} \texttt{p05} --- a problem GPT-OSS
solves alone: routing GPT's working proposal through Qwen for repair destroyed a
capable trajectory. And in the integrated agent's five wins, \emph{none} came from
a handoff rung.

This is a clean negative result, and I think it is the most useful thing in
Part~Two. The mechanism is intuitive: a second model does not share the first's
plan, so it tends to ``repair'' by rewriting toward its own approach, discarding a
trajectory that was nearly closed. Same-model repair preserves the plan; cross-
model repair breaks it. The design consequence is wired into the ladder: run
same-model champions first, and keep cross-model handoff only as a late, seeded
step --- never as a default.

\subsection{Failure taxonomy}
Classifying the \Sdev{} failures into the task's four categories
(Appendix~\ref{app:failtax}) shows the bottleneck is formalization and stuck
repair, not incorrect informal reasoning or harness errors: \texttt{p09} and
\texttt{rmo\_2000\_2} are stuck-repair near-misses (both are reachable in
principle), \texttt{rmo\_2000\_3} fails formalization outright, and
\texttt{putnam\_2018\_a1} produces only guard-rejected tautologies.

\subsection{Verification is part of the science}
The excluded Putnam row is not a footnote. Without a statement lock, a model earns
a comparator ``pass'' by redefining the answer set to be the question itself
(\S\ref{sec:mechanism}); the tautology lands on the GPT-side arms, so counting it
would inflate both parts of this study and even distort Finding~1. Guarding against
it is what makes the rest of the numbers trustworthy.
